\section{Stationarization from recurrent narrow cuts}\label{sec:ramsey56}
The first reduction does not require reversible physical dynamics. In this section only, replace $H$ by a topological monoid $G$ with identity $1$, retain a finite command alphabet, and let $F_{s,j}:G\to Y_j$ be continuous. The output sets are norm-compact convex subsets of real Banach spaces. Define $M_\varepsilon$ and $W_{N,\varepsilon}$ by the same stochastic models. Compactness of the physical monoid is not required.

\begin{definition}[Eventual approximate returns]\label{def:returns56}
The alphabet has eventual approximate returns if, for every neighborhood $V$ of $1$, there is an integer $g(V)\ge0$ such that every integer $n\ge g(V)$ is the length of a word with product in $V$. Cofinite exact returns mean that a single $g$ works with product exactly $1$.
\end{definition}
The statement concerns physical products only. The hidden rows on a return word need not be close to an identity matrix. Approximation here is eventual at every length, not merely along a return subsequence.

\begin{theorem}[Width-preserving stationarization]\label{thm:stationarization56}
Assume eventual approximate returns. For every $\varepsilon\ge0$ and integer $k\ge1$, the following are equivalent:
\begin{enumerate}
\item a stationary realization with at most $k$ labels has error at most $\varepsilon$;
\item among error-$\varepsilon$ clocked realizations, the number of positive cuts having width at most $k$ is unbounded.
\end{enumerate}
Consequently, if $M_\varepsilon>k$, there is $L_{k,\varepsilon}<\infty$, independent of the horizon, such that every error-$\varepsilon$ realization satisfies
\begin{equation}\label{eq:occupation56}
 \#\{t\in\{1,\ldots,N\}:|S_t|\le k\}\le L_{k,\varepsilon}.
\end{equation}
All other registers may be arbitrarily large. Moreover,
\begin{equation}\label{eq:limitminimum56}
                  \lim_{N\to\infty}W_{N,\varepsilon}=M_\varepsilon
                  \quad\hbox{in }\N\cup\{\infty\}.
\end{equation}
If the right side is finite, equality holds for every sufficiently large $N$.
\end{theorem}
The theorem preserves the proposed width, including randomized initialization. It does not assert that the microscopic rows of the clocked machine become stationary. They can be entirely unrelated at different times.

\subsection{Finite witnesses and homogeneous transition tables}
For a matrix $A$ write $\|A\|_{\mathrm r}=\max_i\sum_j|A_{ij}|$. Put $a=|\mathcal A|$ and $D_* =\max(1,\max_j\operatorname{diam}Y_j)$. Let $\mathcal X_k$ be the compact parameter space of stationary $k$-label machines. For $m\ge0$ define
\[
 e_{k,m}=\min_{X\in\mathcal X_k}\max_{s,j,\,|w|\le m}
               \|X_{s,j}(w)-F_{s,j}(u_w)\|_j.
\]
The finite maximum is continuous, so the minimum is attained. If no stationary error-$\varepsilon$ machine exists, compactness of $\mathcal X_k$ and the finite intersection property give an $m$ with
\begin{equation}\label{eq:finitewitness56}
                       e_{k,m}=\varepsilon+\gamma,\qquad\gamma>0.
\end{equation}
Indeed the nested closed sets defined by the constraints $|w|\le m$ otherwise have a common point. This argument also applies when the full infinite-word error is not a continuous function of the parameters.

We shall use the finite edge-coloring theorem of Ramsey \cite{Ramsey}. Denote by $R_q(h)$ an integer such that every coloring of the edges of a complete graph on $R_q(h)$ vertices with at most $q$ colors has a monochromatic $h$-vertex subgraph. Only finiteness is essential. For definiteness one may take
\begin{equation}\label{eq:ramseybound56}
                   R_q(h)\le q^{qh+1}\qquad(q\ge2,\ h\ge2).
\end{equation}
To see this elementary bound, successively select a vertex and retain a largest color class among its remaining neighbors. With $q^{qh+1}$ initial vertices this produces at least $q(h-1)+1$ selected vertices, each of whose edges to later selected vertices has one color. At least $h$ of their associated colors agree; these vertices form the required subgraph. The deliberately excessive exponent avoids any endpoint rounding issue.

Partition $[0,1]$ into at most $\lceil k/\delta\rceil+1$ intervals of length at most $\delta/k$. Two tuples in the same entrywise cell of
\[
                  \mathsf{Stoch}(k)^{\{0\}\cup\mathcal A}
\]
have corresponding matrices within $\delta$ in $\|\cdot\|_{\mathrm r}$. The number of occupied cells is at most
\begin{equation}\label{eq:colors56}
             q=\max\bigl(2,(\lceil k/\delta\rceil+1)^{(a+1)k^2}\bigr).
\end{equation}
Each representative used below is chosen from an occupied cell and is therefore stochastic; an entrywise rounded matrix need not be stochastic and is not used as a representative.

\begin{proof}[Proof of Theorem~\ref{thm:stationarization56}]
The forward implication follows by using one stationary machine at every horizon. For the converse, suppose \eqref{eq:finitewitness56}. Choose $0<\eta<\gamma/4$. Continuity of multiplication and of the finitely many $F_{s,j}$ gives a neighborhood $V$ of $1$ with the following property. In any word of length at most $m$, inserting at most $m+3$ factors from $V$, including before and after the word, changes every output by less than $\eta$. There are only finitely many relevant command words and insertion patterns; continuity at the tuple of identity factors, followed by a finite intersection of neighborhoods, proves the assertion. This uses neither an invariant metric nor commutativity.

Let $g=g(V)$, put $b=g+1$, and choose
\[
       0<\delta\le \min\bigl(1,\gamma/(4D_*(m+1))\bigr).
\]
Take a clocked realization with so many narrow cuts that, after discarding cuts $t<g$ or $t>N-g$ and greedily retaining cuts separated by at least $b$, at least $R_q(m+2)$ remain. Such a choice is possible by hypothesis; at most $2g$ cuts are discarded and each retained cut accounts for at most $b$ of the others.

Identify each retained register with a subset of $\{1,\ldots,k\}$. Extend rows from dummy labels arbitrarily to stochastic rows without altering any row on the actual register. For retained cuts $t_i<t_j$ write $\ell=t_j-t_i\ge g+1$. Choose a length-$\ell$ return word with product in $V$, and let $A_{ij}$ be its actual composite kernel between these cuts. For each letter $c\in\mathcal A$, choose the word consisting of $c$ followed by a length-$(\ell-1)$ return word with product in $V$, and let $B_{ij,c}$ be its composite kernel. Thus the physical products of these words are respectively in $V$ and of the form $v u_c$ with $v\in V$. All the matrices are of order $k$, however large the intermediate registers may be.

Color the edge $\{i,j\}$, oriented from the earlier cut to the later cut, by the cell of $(A_{ij},(B_{ij,c})_c)$. Ramsey's theorem gives $m+2$ ordered cuts
\[
                         t_0<t_1<\cdots<t_{m+1}
\]
whose edges have one color. Fix a stochastic representative $(E,(B_c)_c)$ from that color. Choose fixed return words for the prefix of length $t_0$ and the suffix of length $N-t_{m+1}$, both with products in $V$. Let $\alpha_s$ be the actual state distribution after the prefix. Let $D_j(i)$ be the conditional expected original decoder after the suffix, given label $i$ at its start, with arbitrary legal values on dummy labels. Convexity gives $D_j(i)\in Y_j$.

Consider the stationary machine
\begin{equation}\label{eq:extractedmachine56}
                  \alpha_s,\qquad T_c=B_c,\qquad d_j=ED_j.
\end{equation}
Its decoder is legal. To test $w=c_1\cdots c_\ell$, $0\le\ell\le m$, use the first $\ell$ consecutive clique edges with kernels $B_{i-1,i,c_i}$ and then the single edge from $t_\ell$ to $t_{m+1}$ with kernel $A_{\ell,m+1}$. For $\ell=0$ there is only this last edge. Add the fixed prefix and suffix. This is a genuine word of length $N$; its physical product is obtained from $u_w$ by at most $\ell+3$ insertions from $V$. Its target output is consequently within $\eta$ of $F_{s,j}(u_w)$.

Each of its $\ell+1$ edge matrices differs from the corresponding matrix in \eqref{eq:extractedmachine56} by at most $\delta$. Stochastic multiplication contracts the $\ell^1$ distance between probability rows. A telescoping expansion, and subtraction of any fixed point of $Y_j$ from the decoder, therefore bound the difference between the actual output and $\alpha_s B_{c_1}\cdots B_{c_\ell}E D_j$ by
\[
                           D_*(\ell+1)\delta.
\]
Wordwise correctness of the original machine implies $e_{k,m}\le\varepsilon+\eta+D_*(m+1)\delta<\varepsilon+\gamma$, a contradiction. The representative $E$ is used only as part of the terminal decoder. No idempotence of $E$, consistency between different filler words, or hidden-state identity on a physical return was assumed.

The same argument shows that one may take
\begin{equation}\label{eq:ramseyloss56}
              L_{k,\varepsilon}=2g+(g+1)R_q(m+2).
\end{equation}
For exact cofinite returns, use exact fillers and omit $\eta$ and the neighborhood choice; $g$ is the given exact-return bound. In general $g$ depends on the finite witness and the continuity neighborhood. The remaining parameters depend only on the interface, the alphabet, $k$, and $\varepsilon$, never on a machine or horizon.

If $M_\varepsilon=\infty$, \eqref{eq:occupation56} for each $k$ gives $W_{N,\varepsilon}\to\infty$. If $M_\varepsilon=K<\infty$, the stationary machine gives $W_{N,\varepsilon}\le K$ for every $N$, and the bound for $k=K-1$ gives eventual equality when $K>1$. When $K=1$ equality holds for every $N$ because registers are nonempty.
\end{proof}

\begin{proposition}[Approximate returns without exact returns]\label{prop:approxreturns56}
Let the two circle commands have angles $\theta,\phi$, where $1,\theta/(2\pi),\phi/(2\pi)$ are rationally independent. They have no nonempty exact-return word, but satisfy Definition~\ref{def:returns56}.
\end{proposition}
\begin{proof}
A length-$n$ product has angle $n\phi+j(\theta-\phi)$, $0\le j\le n$. The positive orbit of the irrational rotation $\theta-\phi$ is dense. Compactness supplies, for every accuracy, a finite initial orbit segment that is a net of the circle. All longer initial segments remain nets; rotation by $n\phi$ preserves their accuracy. Thus a product is arbitrarily close to the identity at every sufficiently large length. An exact return would be a nontrivial integer relation among $1,\theta/(2\pi),\phi/(2\pi)$.
\end{proof}
A single irrational letter fails the eventual condition: there is only one word at each horizon. Its target can be placed in a one-label horizon-dependent decoder even when no finite stationary approximation at the same tolerance exists. Thus mere recurrence along a subsequence is not a substitute for Definition~\ref{def:returns56}.
